\thispagestyle{empty}
\vspace*{\fill}

{\Large\bfseries Blackship Technical Manual}\par
\vspace{0.5em}
Version \BlackshipVersion, \ManualEdition

\vspace{2em}
Copyright \textcopyright\ 2026 The Blackship contributors.

This manual and the Blackship source code are distributed under the BSD
3-Clause License. The authoritative license text is in the project repository.

The cover uses original raster artwork. Blackship takes its name and atmosphere
from the broad science-fiction idea of ominous ``black ships''; this project and
manual are not affiliated with or endorsed by Games Workshop, and contain no
Warhammer artwork, logos, characters, or setting text.

\vspace{2em}
\begin{shipnote}[Manual scope]
This edition documents Blackship \BlackshipVersion. The command line remains the
authoritative interface: run \command{blackship --help}, a subcommand's help,
or \command{blackship mcp}'s generated manual when using a newer release.
\end{shipnote}

\vspace*{\fill}
\clearpage

\chapter*{How to read this manual}
\addcontentsline{toc}{chapter}{How to read this manual}

Chapters 1--4 orient a new operator and lead to a running jail. Chapters 5--9
form the day-to-day reference for configuration, images, networking, storage,
operations, and commands. Chapters 10--11 cover integration and the security
model. Chapter 12 is organized by symptoms for incident response.

Shell sessions use a dollar sign for an unprivileged prompt and a hash sign for
root. Do not type the prompt character. Examples assume FreeBSD 15.1 unless the
text says otherwise. Replace pool names, interfaces, releases, addresses, and
paths with values appropriate to the host.

Destructive commands are called out explicitly. A dry run is recommended before
any dependency-wide lifecycle or reclamation operation.
\clearpage
